% !TEX root = riemannian_complex_euclidean.tex
\documentclass[11pt]{article}

\usepackage{amsmath}
\usepackage{amssymb}
\usepackage{amsthm}
\usepackage{geometry}
\usepackage{booktabs}
\usepackage{enumitem}
\usepackage{parskip}

\geometry{margin=1in}

\title{Formal Framing: Geometric Routing via Riemannian--Complex Structure}
\author{}
\date{}

\begin{document}

\maketitle

\section{Core Hypothesis}

We propose that standard transformer architectures are constrained by their reliance on
Euclidean linear operators, which inadequately capture directional, phase-dependent, and
hierarchical structure.

We introduce a geometric routing framework in which:
\begin{itemize}
    \item Local dynamics are governed by a complex-valued transport law
    \item Global structure is encoded in a Riemannian manifold
    \item Information propagation follows geodesic-aware routing rather than purely linear mixing
\end{itemize}

\section{Geometric State Space}

Let $(M, g)$ be a smooth Riemannian manifold with metric $g$.
Each token or latent state evolves as:
\[
    x_t \in M
\]

\section{Tangent Space Representation}

At each point $x_t$, we define a local representation in the tangent space:
\[
    v_t \in T_{x_t}M
\]
To introduce directional and phase-aware structure, we extend this to a complexified tangent
space:
\[
    z_t \in T_{x_t}M \otimes \mathbb{C}
\]
with decomposition:
\[
    z_t = v_t^{(r)} + i\,v_t^{(i)}
\]

\section{Local Transport Operator (Complex Phase Dynamics)}

We define a local update operator:
\[
    \mathcal{U}_{x_t} \colon T_{x_t}M \otimes \mathbb{C}
    \;\rightarrow\; T_{x_t}M \otimes \mathbb{C}
\]
which encodes:
\begin{itemize}
    \item rotation (phase evolution)
    \item scaling (amplitude modulation)
    \item directional coupling
\end{itemize}

A minimal form:
\[
    \mathcal{U}_{x_t}(z) = W_r z + i\,W_i z
\]
where $W_r, W_i$ are learned linear maps on $T_{x_t}M$.

This generalizes real-valued linear layers by allowing intrinsic rotational structure rather
than encoding it implicitly.

\section{Manifold Update (Geometric Routing Step)}

The updated state is obtained via exponential mapping:
\[
    x_{t+1} = \exp_{x_t}\!\Big(\Re\bigl(\mathcal{U}_{x_t}(z_t)\bigr)\Big)
\]
Key properties:
\begin{itemize}
    \item respects manifold geometry
    \item enforces constraint structure (e.g., curvature, boundedness)
    \item replaces unconstrained linear accumulation
\end{itemize}

\section{Interpretation}

This framework separates computation into two regimes:

\textbf{Local (Tangent Space)}
\begin{itemize}
    \item complex-valued
    \item phase-aware
    \item linearizable
    \item supports rotation/interference
\end{itemize}

\textbf{Global (Manifold)}
\begin{itemize}
    \item nonlinear
    \item curvature-aware
    \item enforces structure
    \item governs routing and compression
\end{itemize}

\section{Relationship to Standard Architectures}

\begin{center}
\begin{tabular}{lll}
\toprule
\textbf{Component} & \textbf{Standard Transformer} & \textbf{Proposed Framework} \\
\midrule
State space  & $\mathbb{R}^n$                   & Riemannian manifold $M$       \\
Local ops    & real linear layers               & complex transport operators   \\
Mixing       & attention / linear projection    & geodesic-aware routing        \\
Geometry     & implicit / flat                  & explicit / curved             \\
\bottomrule
\end{tabular}
\end{center}

\section{Key Claim}

The advantage does not arise from complex numbers alone.
Rather, it emerges from:

\begin{quote}
The interaction between complex local transport (phase dynamics) and global Riemannian
structure (geodesic routing), which induces non-linear compression and structured information
flow not achievable in purely Euclidean linear systems.
\end{quote}

\section{Minimal Theoretical Statement}

Let $M$ be a Riemannian manifold and $T_xM \otimes \mathbb{C}$ its complexified tangent
bundle. A routing process defined by local complex-linear transport followed by exponential
mapping induces a nonlinear flow on $M$ whose trajectories encode both phase-dependent
dynamics and curvature-constrained propagation.

\section{Intuition}

\begin{itemize}
    \item \textbf{Euclidean models:} information spreads like diffusion in a flat grid.
    \item \textbf{Proposed model:} information moves like waves constrained to a curved
          surface, with phase, interference, preferred directions, and natural compression
          into geometric sectors.
\end{itemize}

\section{What \textquotedblleft$2i$\textquotedblright\ Actually Becomes}

\begin{itemize}
    \item $i \;\to\;$ local $90^{\circ}$ rotation operator
    \item $2i \;\to\;$ rotation $+$ scaling
\end{itemize}

In this framework, scalar complex multiplication generalizes to a learned complex transport
field over a manifold:
\[
    \text{scalar } 2i
    \;\;\longrightarrow\;\;
    \mathcal{U}_{x_t}(z) = W_r z + i\,W_i z
\]

\end{document}
